problem:  
Let \((X,d,\mathfrak m)\) be an \(RCD(K,N)\) metric measure space with \(K\in\mathbb R\) and finite \(N>1\).  
For Lipschitz \(f:X\to\mathbb R\), denote by \(u_t=P_t f\) the gradient flow of the Cheeger energy (the heat semigroup) and write  
\[
E_t(f):=\int_X |\nabla u_t|^2\,d\mathfrak m,\qquad E_0(f):=\int_X |\nabla f|^2\,d\mathfrak m.
\]  
It is known that \(E_t(f)\) is nonincreasing and that \(E_t(f)\le e^{-2Kt}E_0(f)\).

Define the **normalized energy ratio**:
\[
\mathcal R_t(f):=e^{2Kt}\frac{E_t(f)}{E_0(f)}=1-r(f)\,t^2+o(t^2),\quad t\to0^+.
\]
In smooth Bakry–Émery manifolds, \(r(f)\) encodes both \(\operatorname{Ric}_V\) and its derivatives.  

**Problem:**  
Is there an intrinsic notion of the second-order curvature functional \(r(f)\) in the nonsmooth \(RCD(K,N)\) framework whose spatial constancy characterizes spaces of *parallel synthetic curvature*, i.e., whose tangent cones are all Gaussian after rescaling?  
Equivalently, does small-time *quadratic flatness* of the normalized energy decay, uniform across test functions, imply that \((X,d,\mathfrak m)\) is isometric (up to scaling) to a weighted Euclidean space \((\mathbb R^n,|\cdot|,e^{-\frac K2|x|^2}dx)\)?  

Exploring further, can one define \(r(x)\) as the limit of suitably normalized second differences in \(E_t(f)\) for sharply localized \(f\) near \(x\), and, if so, does constancy of \(r(x)\) entail constancy of the “synthetic Hessian of the entropy” or of the metric–measure Laplacian spectrum?

Why is it a "good" problem:  
This problem pushes curvature–dimension analysis beyond first-order Bochner-type identities into second-order asymptotics, a regime not yet defined in general for \(RCD(K,N)\) spaces. It invites the construction of a *second-order curvature invariant* based on the fine decay of the heat-flow energy ratio—connecting probabilistic diffusion asymptotics, synthetic Ricci geometry, and spectral rigidity.  
Unlike first-order expansions, the second-order term involves curvature derivatives and genuinely new geometric content; defining and interpreting it in the nonsmooth synthetic setting would demand fundamentally new techniques.  
A resolution would inaugurate a “second-order” Bakry–Émery theory on metric measure spaces, revealing whether uniform higher-order diffusion symmetry forces Gaussian–type rigidity. This clean yet profound question thus lies at the frontier where analytic expansions meet metric geometric structure.